\chapter{Parameter Definitions and Conventions}
\label{ch:params}
% Preserve the historical community reference without treating it as primary data.
\nocite{perfectgolfswing}

A launch-monitor number describes a chosen quantity, at a chosen point and time, through a particular estimation procedure.
A disagreement can therefore reflect different definitions as well as measurement error.
This chapter uses TrackMan's published conventions where stated, without treating them as a universal standard or an independent accuracy validation~\cite{trackman40params,trackmanclubdata,trackmanparameterdefinitions}.
The right-handed world frame used here and in \cref{ch:impact} has $x$ forward along the target line, $y$ upward and $z$ right.
Horizontal angles are positive right of target and elevations are positive upward; right and left do not change with the golfer's handedness.
Terms such as open, closed, in-to-out and draw need an additional handedness convention.

\section{Club-Delivery Parameters}
\label{sec:clubparams}

The time assigned to an output need not be the time interval used to estimate it.
TrackMan's 2017 account describes club-motion estimates fitted from pre-impact observations and referenced to maximum compression, excluding collision-induced motion~\cite{trackmanclubdata}.
Its current parameter list explicitly defines club speed just before first contact, while path and attack angle retain the maximum-compression reference~\cite{trackmanparameterdefinitions,trackmanclubspeed}.
Thus the latter reference does not justify interpreting every displayed club parameter as an instantaneous measurement of the colliding head.
Keep the device, software version, output definition and observation window with a comparison dataset.

\begin{table}[htbp]
\centering\small
\caption{Club-Delivery Quantities and Their Declared Conventions}
\label{tab:clubparams}
\begin{tabular}{@{}p{3.0cm}p{7.0cm}p{4.6cm}@{}}
\toprule
\textbf{Parameter} & \textbf{Quantity} & \textbf{Interpretation} \\
\midrule
Club speed & Norm of geometric-center (GC) velocity before first contact & Reference-point dependent \\
Attack angle & Elevation of GC motion at the stated reference event & Positive upward \\
Club path & Azimuth of GC motion at that event & Positive right of target \\
Face angle & Azimuth of the local face normal at the contact center & Positive right of target \\
Face to path & Face angle minus club path & A difference of headings \\
Dynamic loft & Elevation of the local contact normal & Includes delivered pose and face curvature \\
Spin loft & Angle between the declared motion direction and face normal & A three-dimensional angle \\
Swing plane / direction & Inclination / horizontal orientation of a plane fitted to a clubhead trajectory segment & Declare the fitted interval and point \\
Low point & Location of the swing-arc minimum relative to the impact reference & Does not alone determine turf contact \\
Impact height / offset & Strike coordinates relative to face center & Declare face axes and signs \\
Dynamic lie & TrackMan: lower-shaft inclination to horizontal at impact & Not identical to face roll \\
Closure rate & Foresight: heel-to-toe rotation rate about the shaft & Not the full angular-velocity norm \\
\bottomrule
\end{tabular}
\end{table}

The delivery definitions in \cref{tab:clubparams} follow the cited vendor documentation; Foresight's lie output instead describes face-plane roll~\cite{trackmanparameterdefinitions,foresightclubdefinitions}.
A shaft inclination and a face-roll coordinate cannot be compared by name alone.
Likewise, a rate about a shaft axis does not supply the other components of a clubhead's angular velocity, and a finite-window rate need not equal its value at one instant.
No exclusivity claim about which product can report these quantities follows from their definitions.

A fitted swing plane summarizes a selected trajectory segment; it does not imply that the whole swing is planar.
TrackMan's low-point convention locates the arc minimum relative to the GC at maximum compression~\cite{trackmanlowpoint}.
An arc minimum ahead of that reference does not establish ball-first contact with the ground: absolute height, the sole and leading-edge geometry, ball position and the ground surface also matter.
A translated copy of the same arc can have the same fore-aft low point while intersecting the ground at a different place.

Face angle and dynamic loft refer here to the normal at the contact center.
Bulge and roll make that normal differ from the face-center normal on an off-center driver strike (\cref{sec:gear}).
A measurement of face-center pose is not intrinsically wrong; it answers a different question unless local surface geometry and strike location are supplied.
Similarly, face-to-path can help interpret a centered strike under a stated collision model, but its sign alone does not determine curvature in the presence of gear effect or differing flight conditions.

\subsection{Reference Point and Rigid-Body Velocity}
\label{sec:pathreference}

For two material points on an ideal rigid clubhead at the same instant,
\begin{equation}
\label{eq:pointvel}
\vect{v}_P=\vect{v}_O+\boldsymbol\omega\times\vect r,
\qquad \vect r=\vect p_P-\vect p_O.
\end{equation}
All vectors must be expressed in the same frame.
Both speed and direction can change with the reference point; geometric center, center of mass and face center are distinct definitions even when two happen to be close.
TrackMan's 2017 vendor illustration reports an approximately $3\degs$ driver path difference between center of gravity and face center, with the face-center path farther left~\cite{trackmanclubdata}.
That illustration is neither a universal correction nor evidence that every radar and optical product chooses those respective points.

For a manufactured example in SI units, choose $\vect v_O=(40,0,0)$, $\vect r=(0.04,0,0)$ and $\boldsymbol\omega=(0,50,0)$.
Then $\vect v_P=(40,0,-2)$ and its horizontal heading is $\operatorname{atan2}(-2,40)\simeq-2.86\degs$.
Reversing $\boldsymbol\omega$ reverses the heading offset; choosing $\boldsymbol\omega=(50,0,0)$ gives zero offset despite the same angular-speed norm.
With horizontal speed fixed and nonzero, increasing upward velocity increases attack angle; for an already upward motion, that makes the path steeper.
These examples demonstrate geometry, not typical golfer values or measured device discrepancies.

\Needspace{10\baselineskip}
\begin{keypoint}
A rigid-body twist $(\boldsymbol\omega,\vect v_O)$ and known $\vect r$ permit conversion between reference points (\cref{app:screw}).
The scalar $\|\boldsymbol\omega\|$ alone does not.
The velocity difference obeys $\|\Delta\vect v\|\leq\|\boldsymbol\omega\|\|\vect r\|$, but converting that bound to a heading difference also requires the original velocity and a nonvanishing horizontal projection.
Two devices can be compared only after reconciling their reference points, frames and times, with the conversion uncertainty retained.
\end{keypoint}

\subsection{Spin Loft and Contact Kinematics}

Write attack angle as $a$, club path as $p$, dynamic loft as $\ell$ and face angle as $f$.
For nonvertical directions, the unit motion vector is $(\cos a\cos p,\sin a,\cos a\sin p)$ and the unit face normal has the analogous form with $(\ell,f)$.
Their dot product gives the spin-loft angle $\sigma\in[0,\pi]$:
\begin{equation}
\label{eq:parametersspinloft}
\cos\sigma=\sin a\sin\ell+\cos a\cos\ell\cos(f-p).
\end{equation}
This implements the vendor's three-dimensional angular definition~\cite{trackmanspinloft}.
When $f=p$ and the elevation difference lies in $[-\pi,\pi]$, $\sigma=|\ell-a|$; using a signed difference silently assumes its sign.
For $a=-5\degs$, $\ell=15\degs$, changing $f-p$ from zero to $10\degs$ increases spin loft from $20\degs$ to approximately $22.32\degs$.
Direct vector evaluation still works when an azimuth is undefined at a vertical direction, provided both vectors are known.

Reported face-to-path and spin loft are legitimate declared quantities even when face orientation is evaluated at contact and velocity at GC.
They are not automatically the local relative contact kinematics required by an impact law.
For ideal rigid club and ball kinematics immediately before contact, the club-point velocity relative to the prospective ball contact point is
\begin{equation}
\vect v_{\rm rel}=\vect v_P-
 \left[\vect v_B+\boldsymbol\omega_B\times(\vect p_Q-\vect p_B)\right],
\end{equation}
where $P$ and $Q$ are the respective surface points and $B$ is the ball center.
During a deformable impact, local surface-deformation velocities must also be included; the displayed rigid-body expression alone does not represent them.
Its normal and tangential components enter a contact model; compliance, friction and impact location then influence the outgoing ball state (\cref{ch:impact}).
Knowing two headings does not identify all of those quantities.

\section{Ball-Launch Parameters}

\begin{table}[htbp]
\centering\small
\caption{Launch State and Separately Derived Flight Descriptors}
\label{tab:ballparams}
\begin{tabular}{@{}p{3.2cm}p{10.8cm}@{}}
\toprule
\textbf{Parameter} & \textbf{Meaning and Scope} \\
\midrule
Ball speed & Norm of center velocity immediately after separation \\
Launch angle / direction & Elevation / azimuth of that outgoing velocity \\
Spin rate & Angular-speed magnitude, usually expressed in rpm \\
Spin axis & Direction of the nonzero spin vector; one tilt angle requires an additional convention \\
Smash factor & Ball-speed / club-speed ratio, using their respective events and reference points \\
Carry / carry side & Distance / lateral position at a declared flight endpoint, not an instantaneous launch state \\
Total & Distance to a resting position, requiring bounce and roll treatment \\
Apex / landing angle & Maximum flight height / velocity elevation at a declared landing event \\
\bottomrule
\end{tabular}
\end{table}

TrackMan's carry convention uses the descending crossing of the launch elevation, not necessarily the first contact with actual terrain~\cite{trackmancarry}.
Its parameter documentation distinguishes this flight descriptor from initial launch quantities and describes total using a calculated resting position~\cite{trackmanparameterdefinitions}.
A launch-only indoor measurement therefore cannot by itself observe carry, landing or roll.
Comparisons must state whether flight is tracked, extrapolated or normalized to chosen environmental conditions (\cref{ch:flight}).
Smash factor is a speed ratio, not a fraction of kinetic energy transferred: it omits the relevant masses and rotational and deformational energy.

\subsection{Spin Components Need a Frame}

An ideal rigid ball has one instantaneous angular-velocity vector, not independent physical ``backspin'' and ``sidespin'' rotations.
Use the orthonormal launch frame in \cref{eq:axistilt}: outgoing direction $\uvec v_b$, transverse backspin axis $\uvec e_b$ and transverse up axis $\uvec e_u$.
Express the spin vector $\vect S$ in rpm, so $\vect S=60\boldsymbol\omega_b/(2\pi)$ when angular velocity is in radians per second.
Define $S_{\rm back}=\vect S\cdot\uvec e_b$, $S_{\rm side}=-\vect S\cdot\uvec e_u$ and $S_{\parallel}=\vect S\cdot\uvec v_b$.
Then
\begin{equation}
\label{eq:spincomponents}
\begin{aligned}
 S_\perp&=\sqrt{S_{\rm back}^2+S_{\rm side}^2},&
 \theta&=\operatorname{atan2}(S_{\rm side},S_{\rm back}),\\
 S_{\rm side}&=S_\perp\sin\theta,&
 S_{\rm back}&=S_\perp\cos\theta,\\
 S^2&=S_\perp^2+S_\parallel^2,& S&=\|\vect S\|.
\end{aligned}
\end{equation}
Replacing $S_\perp$ by total spin $S$ is valid only when the axial component vanishes.
For horizontal forward flight and $\vect S=(1000,-300,3000)$ rpm in the world frame, $S_\perp\simeq3014.96$ rpm, $S\simeq3176.48$ rpm and $\theta\simeq5.71\degs$.
Thus total spin and one tilt angle do not determine the general three-dimensional state.
The tilt is undefined if both transverse components vanish, and zero spin has no axis; the chosen launch-frame construction also needs a nonvertical launch.
Device-specific sidespin and tilt conventions must be converted explicitly before applying these equations.

In the usual spin-induced lift model, the transverse spin relative to air velocity contributes to lateral and vertical lift; wind can make that velocity differ from ground-relative launch velocity.
A fixed lateral-distance-per-degree rule cannot replace the flight calculation: speed, launch direction, spin magnitude, wind and aerodynamic coefficients also matter (\cref{ch:flight}).
Left/right curvature labels need not be draw/fade labels for both handednesses.

\section{The Directness Hierarchy}
\label{sec:hierarchy}

Here, \emph{measured} means estimated from observations that constrain the quantity, with calibration and a stated observation model; it does not mean algorithm-free or error-free.
\emph{Derived} denotes a transformation of reported quantities, and \emph{estimated from a physical model} denotes an inference or prediction that additionally depends on collision, flight or ground-interaction assumptions.
These are provenance descriptions rather than mutually exclusive product classes.
Radar records signals and cameras record images; neither records a calibrated three-dimensional club path as raw data.
\Cref{tab:hierarchy} therefore lists the evidence needed for an output rather than assigning every device of one architecture an identical capability.

\begin{table}[htbp]
\centering\small
\caption{Output Provenance and Conditions to Check}
\label{tab:hierarchy}
\begin{tabular}{@{}p{3cm}p{6cm}p{5.5cm}@{}}
\toprule
\textbf{Output} & \textbf{Possible Provenance} & \textbf{Necessary Qualification} \\
\midrule
Ball speed / launch & Measured through calibrated radar or image observations & Geometry, timing, track association and launch-event fit \\
Spin rate / axis & Measured from rotation-sensitive observations; or model-inferred from other quantities & Visible features or spin signatures, ambiguity, axial component and flight window \\
Club motion & Measured by tracking a declared point or fitting rigid motion & Reference point, observation interval and collision exclusion \\
Face pose / strike & Measured from resolved geometry; or inferred through an impact model & Surface normal, strike location, markers if required, and model identifiability \\
Face to path / smash & Derived from the constituent outputs & Shared uncertainty and definition compatibility \\
Carry / landing & Derived from a sufficiently observed trajectory, or predicted from launch & Observed endpoint versus extrapolation and normalization \\
Total & Estimated with flight, bounce and roll models when rest is not observed & Ground slope, restitution, friction and endpoint definition \\
\bottomrule
\end{tabular}
\end{table}

The radar and camera chapters develop these conditional observation relationships (\cref{ch:radar,ch:camera}).
An outdoor radar track can support trajectory-based spin inference, while a short indoor track may require different information; neither statement identifies every product's algorithm.
For example, TrackMan now describes iO as using cameras for full three-dimensional spin, including axial spin~\cite{trackmaniospin}.
That product-specific vendor claim is evidence against a universal ``indoor radar means estimated axis'' label, not independent proof of its advertised accuracy.
Similarly, a marker requirement belongs to the particular club-tracking implementation, not to optical sensing as a physical necessity.

\begin{implication}
To use launch-monitor data to investigate the golf swing, retain the measurement chain: club motion and contact geometry constrain impact; impact produces launch; launch and environment determine flight.
A face angle inferred from ball launch is not independent corroboration of a proposed relationship between those same outputs.
Report definitions, uncertainty and shared inputs alongside the numbers before attributing a difference to technique or to one sensor family.
\end{implication}
